%% PhD thesis template
%% School of Computing and Communications, Lancaster University, 

%% Copyright 2020 Andrew Moore and Alistair Baron
%
% This work may be distributed and/or modified under the
% conditions of the LaTeX Project Public License, either version 1.3
% of this license or (at your option) any later version.
% The latest version of this license is in
%   http://www.latex-project.org/lppl.txt
% and version 1.3 or later is part of all distributions of LaTeX
% version 2005/12/01 or later.
%
% This work has the LPPL maintenance status `maintained'.
% 
% The Current Maintainer of this work is Alistair Baron
%
% This work consists of the files main.tex, and the other example tex files referenced within.


%
% Heavily based on the formatting and content of Alistair Baron's thesis https://eprints.lancs.ac.uk/id/eprint/84887/1/2011Baronphd.pdf, Kelly Widdicks's thesis https://eprints.lancs.ac.uk/id/eprint/143606/1/2019widdicksphd.pdf, and Andrew Moore's thesis (not yet publicly available but when is will be found here on his google scholar https://scholar.google.co.uk/citations?user=mJRN_SIAAAAJ&hl=en). The thesis is also somewhat based on the Cambridge thesis template: https://github.com/kks32/phd-thesis-template
%
% This was created by Andrew Moore and Alistair Baron with help from Paul Rayson, and Edward Dearden.
%
% Conforms with Lancaster University MANUAL OF ACADEMIC REGULATIONS AND PROCEDURES (MARP) 2019-20 for POSTGRADUATE RESEARCH. This can be found here: https://www.lancaster.ac.uk/media/lancaster-university/content-assets/documents/student-based-services/asq/marp/PGR-Regs.pdf
% For layout information the most relevant pages of that document are pages 30 and 31 (Appendix 2)

% Please check that your thesis complies with the current regulations. If not available from the above link, details will be on the PGR Moodle space (https://modules.lancaster.ac.uk/course/view.php?id=22985)

% If you have any useful updates for this template, which you think will benefit others, please make these available via github (https://github.com/InfoLab21/scc-thesis-template).


\documentclass[oneside,12pt]{report}
\usepackage[paper=a4paper]{geometry}
\usepackage[utf8]{inputenc}
\usepackage{graphicx}
\usepackage{svg}
\usepackage{hyperref}
\hypersetup{
    colorlinks=false,
    breaklinks=true
}
% All sections and sub sections have to be labelled in the table of contents and thesis. This does this down to the level of subsubsection.
\setcounter{secnumdepth}{4}
\setcounter{tocdepth}{4}

%discourages hyphenated words at end of lines
\hyphenpenalty=5000
\tolerance=1000


% biblography formatting
% NOTE: At the moment both URL and DOI will be displayed, if you have a bib entry that has both URL and DOI please remove the URL so that only one link is displayed, in this case to the DOI. 
\usepackage[style=authoryear,backend=biber,url=true, doi=true, sorting=nyt, natbib=true]{biblatex}
\usepackage{bibentry}
\addbibresource{ref.bib}

% Appendix formatting
\usepackage[title, titletoc]{appendix}
\renewcommand{\appendixpagename}{Appendix}


% for the title page
\usepackage{datetime}
\newdateformat{monthyeardate}{%
  \monthname[\THEMONTH], \THEYEAR}
% For the declaration
% NOTE IF NOT USING OVERLEAF PLEASE REMOVE THIS AND ADD THE WORD COUNT MANUALLY. The software overleaf uses is this: https://app.uio.no/ifi/texcount/
% This has come from overleaf and has been slightly adapted. This counts everything including footnotes and captions I believe and the appendix, but not bibliography.
% https://www.overleaf.com/learn/how-to/Is_there_a_way_to_run_a_word_count_that_doesn%27t_include_LaTeX_commands%3F
\newcommand{\quickwordcount}[1]{%
  \immediate\write18{texcount -0 -sum -merge  #1.tex > #1-words.sum }%
  \input{#1-words.sum}%
}
%
% TO BE FILLED IN
%
\newcommand\thesistitle{Title of your thesis}
\newcommand\authorname{First (Middle) Last name} % Full name
\newcommand\authordegrees{BSc (Hons), MRes} % This does not have to be filled in, but the regulations state that you should include any degrees you have.

\newgeometry{left=38mm, right=25mm, top=25mm, bottom=25mm, includeheadfoot} % This has to come before the header and footer information

% HEADER AND FOOTER STYLE
% https://en.wikibooks.org/wiki/LaTeX/Customizing_Page_Headers_and_Footers
\usepackage{fancyhdr}
\setlength{\headheight}{15pt}
% For preamble
\fancypagestyle{preamble}{ %
  \fancyhf{} % remove everything
  \fancyfoot[C]{\thepage}
  \renewcommand{\headrulewidth}{0pt} % remove lines as well
  \renewcommand{\footrulewidth}{0pt}
}
\pagestyle{preamble}

% For the main pages
\fancypagestyle{main}{
  \fancyhf{}
  %\fancyhead[R]{\thepage}
  \fancyfoot[C]{\thepage}
  \fancyhead[L]{\textit{ \nouppercase{\leftmark}} }
  \fancyhead[R]{\textit{ \nouppercase{\rightmark}} }
  \renewcommand{\headrulewidth}{1pt} % remove lines as well
  \renewcommand{\footrulewidth}{0pt}
}

% For title and chapter pages
\fancypagestyle{plain}{ %
  \fancyhf{} % remove everything
  \fancyfoot[C]{\thepage} %should keep page number
  \renewcommand{\headrulewidth}{0pt} % remove lines as well
  \renewcommand{\footrulewidth}{0pt}
}

\usepackage{lipsum} % just to add random text as an example

%
%


\begin{document}
\pagenumbering{roman}



\begin{titlepage}

\center

\includesvg[scale=0.6]{figures/lu-logo.svg}\vspace{1cm} % This logo was taken from the main Lancaster University website https://www.lancaster.ac.uk/ The same logo can be found, that is not publicly available, in EPS format: https://www.lancaster.ac.uk/current-staff/communications-and-marketing/marketing/resources/#d.en.417655

\huge  \textbf{\thesistitle}

\vspace{2cm}

\Large \textbf{\authorname, \authordegrees}

School of Computing and Communications

Lancaster University

\vfill

\large

A thesis submitted for the degree of

\textit{Doctor of Philosophy}

\vspace{0.5cm}

\monthyeardate\today

\end{titlepage}

I dedicate this thesis to someone. \textbf{(this is not required when submitting your thesis before your viva and you can add a dedication in your final thesis after your viva if you wish.)}

\clearpage

\section*{\centering Declaration}
I declare that the work presented in this thesis is, to the best of my knowledge and belief, original and my own work. The material has not been submitted, either in whole or in part, for a degree at this, or any other university. This thesis does not exceed the maximum permitted word length of 80,000 words including appendices and footnotes, but excluding the bibliography. A rough estimate of the word count is: \quickwordcount{main}\vspace{48pt} 

\noindent \authorname

\clearpage

\begin{center}
\textbf{\thesistitle}

\authorname, \authordegrees.

School of Computing and Communications, Lancaster University

A thesis submitted for the degree of \textit{Doctor of Philosophy}. \monthyeardate\today
\end{center}

\section*{\centering Abstract}

\textbf{This is the beginning of the abstract that according to the regulations should not be any longer than 300 words. See point 14 in appendix 2 of the regulations \url{https://bit.ly/2Q4H43I}.}

\clearpage

\section*{\centering Publications}
(This is an example taken from Andrew Moore's thesis, adapt this to suit your own thesis.)

\textbf{NOTE: If you have done joint research e.g. whereby you are not the first author and you feel you need to state the contributions you have made in those papers that relate to your thesis then create another section like `Contribution Statements' whereby you state your contributions in those papers that relate to the thesis. See Kelly Widdicks thesis (page 11) for a great example. If you do create a `Contribution Statements' section that you may want to revise the your `Declaration', of which Kelly Widdicks declaration is a good example. Kelly's thesis link \url{https://eprints.lancs.ac.uk/id/eprint/143606/1/2019widdicksphd.pdf}. Also for more details on the regulations by the university on declaring contribution statements see point 18 in appendix 2 of \url{https://bit.ly/2Q4H43I}}

\textbf{ADDITIONAL NOTE: If you have any papers under review perhaps speak with your supervisor before adding them here for anonymity reasons.}\vspace{1cm}

Only one publication, shown below, has been created directly from the thesis, from which large portions of this published work is used within chapter 3:\bigskip

\noindent \fullcite{moore-rayson-2018-bringing}\vspace{24pt}

\noindent The following publication have been generated while developing this thesis, and to an extent has guided the thesis into what it has become:\bigskip

\noindent \fullcite{moss-etal-2019-fiesta}\bigskip

\noindent \fullcite{moore-rayson-2017-lancaster}\bigskip

\noindent \fullcite{el-haj-etal-2016-learning}\bigskip

\noindent \fullcite{moore2016domain}

\clearpage

\section*{\centering Acknowledgements}

Acknowledgements you may want to make. \textbf{(this is not required when submitting your thesis before your viva and you can add a dedication in your final thesis after your viva if you wish.)}

\clearpage

\tableofcontents

\clearpage

\listoffigures

\clearpage

\listoftables

\clearpage
\pagestyle{main}
\pagenumbering{arabic}

\chapter{Introduction}
\section{Section Name}
Start of your introduction, with a reference to the relevant appendix \ref{appendix_introduction}. Example of different citation styles \citep{moore-rayson-2018-bringing}, and \citet{moore-rayson-2018-bringing}. Example of a table \ref{table:introduction}:

\begin{table}[!h]
    \centering
    \begin{tabular}{c|c}
        1 & 2 \\
        3 & 4
    \end{tabular}
    \caption{A table}
    \label{table:introduction}
\end{table}

\lipsum[2-10]



\chapter{Middle chapters}
Middle.

\chapter{Conclusions}
\section{Introduction}
Example of a figure \ref{fig:conclusion}:
\begin{figure}[!h]
    \centering
    \includegraphics[scale=0.8]{figures/conclusion/ucrel_logo.png}
    \caption{UCREL logo.}
    \label{fig:conclusion}
\end{figure}

\begin{appendices}

\chapter{Introduction}
\label{appendix_introduction}
\section{Additional Tables}
Some extra tables \ref{table:appendix_introduction} that should go in the appendix.

\begin{table}[!h]
    \centering
    \begin{tabular}{c|c}
        5 & 6 \\
        7 & 8
    \end{tabular}
    \caption{Caption}
    \label{table:appendix_introduction}
\end{table}

\end{appendices}

\clearpage % Required to make sure the references table of contents page number is correct.
\printbibliography[heading=bibintoc,title=References]

\end{document}

